\section{Sketch-Simulation Details}
\label{app:sketch-simulations}

This appendix expands the compact sketch evidence reported in
Section~\ref{sec:v7-tree-comparisons}. The main text uses three anchors:
exact HLL state parity, a broad family summary, and a learned-state
capacity diagnostic. The generated artifacts behind those anchors are kept
here so the main paper can stay readable while the full comparison remains
auditable.

\paragraph{Run provenance.}
The paper figures and tables are copied from the generated sketch-comparison
bundle recorded in
\texttt{classical\_sketches\_figure\_manifest.json}. The copied paper
artifacts live under
\texttt{paper/ctreepo/assets/sketches/}. The figure manifest records leaf
sizes \(16,32,64,128,256\), one random seed, \(145\) aggregate rows, and
\(40\) learned rows.
Official rows use library or analytic sketch implementations; learned rows
are empirical comparisons against the chosen oracle or reference behavior.

\paragraph{Generated grid.}
Table~\ref{tab:app-classical-sketches-compact} is the generated compact
summary. The full row-level grid is stored as
\texttt{assets/sketches/tables/classical\_sketches\_grid.tex}; the
machine-readable version with run metadata is
\texttt{assets/sketches/tables/classical\_sketches\_aggregate.csv}. We keep
the full grid as an artifact rather than typesetting every row in the
manuscript because the grid is a provenance object, while the compact table
and panels below are the readable scientific summary.

\begin{table}[p]
\centering
\scriptsize
\resizebox{\linewidth}{!}{% Auto-generated by treepo.bench.reports.classical_sketches; do not edit.
\begin{tabular}{lllrrrrrrrr}
\toprule
family & query & best official & official & learned $f$ & $f$ RMSE & learned $g$ & $g$ RMSE & learned joint (best) & joint variant & joint RMSE \\
\midrule
distinct & cardinality & theta\_datasketches & 0 & -- & -- & learned\_g\_exact\_distinct\_union\_state\_space & 0 & learned\_joint\_exact\_distinct\_union\_state\_space & fg & 0 \\
frequency & focus\_frequency & -- & -- & -- & -- & learned\_g\_count\_min\_state\_space & 0 & learned\_joint\_count\_min\_state\_space & fg & 0 \\
frequency & top5\_point\_frequency & count\_min\_datasketches & 0 & -- & -- & -- & -- & -- & -- & -- \\
quantile & rank\_at\_q0.5 & quantiles\_floats\_datasketches & 0.002583 & -- & -- & -- & -- & learned\_joint\_quantiles\_reference\_q0.5 & fg & 0.02422 \\
quantile & rank\_at\_q0.95 & tdigest\_double\_datasketches & 0.002462 & -- & -- & -- & -- & learned\_joint\_kll\_reference\_q0.95 & fg & 0.007558 \\
sampling & accumulator\_summary\_sum & tuple\_accumulator\_datasketches & 0 & -- & -- & -- & -- & learned\_joint\_tuple\_summary\_sum\_reference & fg & 0.007204 \\
sampling & total\_weight & varopt\_strings\_datasketches & 0 & -- & -- & learned\_g\_exact\_total\_weight\_state\_space & 0 & learned\_joint\_exact\_total\_weight\_state\_space & fg & 0 \\
set & a\_not\_b & theta\_datasketches & 0 & -- & -- & -- & -- & learned\_joint\_exact\_set\_a\_not\_b & fg & 0.09717 \\
set & intersection & theta\_datasketches & 0 & -- & -- & -- & -- & learned\_joint\_exact\_set\_intersection & fg & 0.2417 \\
set & union & theta\_datasketches & 0 & -- & -- & -- & -- & learned\_joint\_exact\_set\_union & fg & 0.03052 \\
\bottomrule
\end{tabular}
}
\caption{\textbf{Generated compact sketch grid.} Official rows are
reference sketch reductions or controls. Learned rows are empirical
approximations in the chosen state/query class, not claims of recovering the
library implementation.}
\label{tab:app-classical-sketches-compact}
\end{table}

\begin{figure}[p]
\centering
\includegraphics[width=0.48\linewidth]{classical_sketches_distinct_detail.pdf}\hfill
\includegraphics[width=0.48\linewidth]{classical_sketches_frequency_detail.pdf}
\caption{\textbf{Distinct-count and frequency panels.} The distinct-count
panel includes the HLL/Theta/CPC-style cardinality comparisons and learned
counterparts. The frequency panel shows Count-Min and related
point-frequency comparisons.}
\label{fig:app-sketch-distinct-frequency}
\end{figure}

\begin{figure}[p]
\centering
\includegraphics[width=0.48\linewidth]{classical_sketches_quantile_detail.pdf}\hfill
\includegraphics[width=0.48\linewidth]{classical_sketches_set_detail.pdf}
\caption{\textbf{Quantile and set panels.} Quantile rows compare
rank-preserving sketches and learned approximations. Set rows compare union,
intersection, and difference-style queries.}
\label{fig:app-sketch-quantile-set}
\end{figure}

\begin{figure}[p]
\centering
\includegraphics[width=0.64\linewidth]{classical_sketches_sampling_detail.pdf}
\caption{\textbf{Sampling-style panel.} Weighted and accumulator-style
summaries are treated as empirical references unless their state-level
local laws are supplied separately.}
\label{fig:app-sketch-sampling}
\end{figure}

\paragraph{Interpretation.}
The exact HLL row is the state-level parity check: with the native register
state and pointwise-max merge, the tree reduction and flat reduction agree
as states. Broader official rows are empirical references for familiar
sketch families. Learned rows test whether the selected state/query class
can approximate the reference behavior; remaining error is interpreted as
optimization error only when the target law is representable in that class,
and otherwise as a structural gap.
